\section{Discussion}
\label{sec:discussion}

\subsection{Feature Lifting Is a Distinct Agent Capability}

Feature lifting asks a different question from issue repair. A repair agent
succeeds inside the original repository and can rely on its package layout,
resources, and test harness. A feature-lifting agent must infer which parts of
that environment are semantically necessary and reconstruct them behind a new
package boundary. It must satisfy an output contract after the original
repository is removed. This combination makes dependency closure and
modularization part of correctness rather than implementation style.

The task is also different from code localization. Finding a relevant symbol
can be necessary, but the \texttt{requests\_cache} and configuration examples
show why it is insufficient. Conversely, a behaviorally equivalent rewrite can
pass even if it copies little upstream code. \flb therefore evaluates delivered
behavior and artifact independence, not whether the agent followed one
prescribed extraction strategy.

\subsection{Correctness and Compactness Should Remain Separate}

A broad vendoring solution may preserve behavior but offer little reusable
modularization. A very small package may omit rare behavior. Combining these
dimensions into one scalar would obscure both failure modes and make tradeoffs
difficult to interpret. \flb treats Functional Pass as the eligibility gate
for compactness analysis. \rres is then a descriptive reference-relative
measure, not a proof of minimality or semantic authenticity.

The earlier External-50 expansion demonstrates the importance of this
separation. High pass rates alone suggested that the task was nearly solved,
while pass-conditioned footprints revealed that many solutions were close to
whole-repository copying. Hard-50 and its copy-trap tasks retain functional
difficulty while making broad copying visible.

\subsection{Implications for Agent Design}

The negative method results do not imply that agent workflows are irrelevant.
They suggest that a useful workflow must improve the evidence available for
contract closure rather than merely add procedural steps. Promising directions
include explicit obligation tracking tied to public requirements, generation
of boundary cases from exception and state semantics, provenance-aware
dependency closure, artifact-level completeness checks, and calibrated
uncertainty when a requirement cannot be verified legally.

The Public-feedback ablation suggests a layered view of feedback. Executable
Public tests can expose basic mismatches efficiently. Deeper Hidden behavior
remains a generalization problem under the Main information boundary. A robust
agent must use repository evidence and the public specification to anticipate
unobserved but declared combinations without hard-coding benchmark examples.

\subsection{Scope of the Claims}

Our main scores characterize model backends embedded in a fixed OpenHands
\mainprotocol configuration. They are not measurements of base models
independent of runtime, prompt, context policy, or provider behavior. If no
second-runtime experiment is completed, all model and trajectory claims are
explicitly scoped to this protocol. Likewise, task-factor and semantic-failure
analyses are descriptive and do not establish causal effects of repository
properties.

Harness-Bench argues that executable-agent capability should be reported at the
model--harness configuration level~\cite{yao2026harnessbench}. \flb adopts that
reporting discipline by fixing and attesting its Official Main runtime. Our
scientific axis is nevertheless different: we introduce a task construct and
study capability and failure within it, rather than treating harness variation
as the main independent variable.

